\section{Covariant homogeneous OPE coefficients}
\label{sec:ope-homogeneous}
The finite tests of Section~\ref{sec:ope-prescription} are chosen for exact
compatibility, without imposing rotational symmetry. We first average
their field-valued output, retaining both the input and output rotations.
We then construct homogeneous complements using these averaged tests.
This order ensures that the final coefficients have compatible finite
truncations as well as the required symmetries.

Write $R_n(K)$ for the continuous $O(4)$ action on $E_n$ and $J_n$ for
complex conjugation of fields. Assume, as supplied by the field structure,
that these actions commute with the inclusions, that $J_n$ is an antilinear
involution, and that $J_nR_n(K)=R_n(K)J_n$. Assume moreover that the
mixed correlations are Euclidean covariant and real on separated
configurations: for every finite label list $\boldsymbol v$, every
separated position list $\boldsymbol x$, every $K\in O(4)$ and every
$s\in\mathbb R^4$,
\begin{equation}\label{eq:ope-mixed-covariance}
 S_{\mathrm{mix}}(R(K)\boldsymbol v;K\boldsymbol x)
  =S_{\mathrm{mix}}(\boldsymbol v;\boldsymbol x),\qquad
 S_{\mathrm{mix}}(\boldsymbol v;\boldsymbol x+s)
  =S_{\mathrm{mix}}(\boldsymbol v;\boldsymbol x),\qquad
 S_{\mathrm{mix}}(J\boldsymbol v;\boldsymbol x)
  =\overline{S_{\mathrm{mix}}(\boldsymbol v;\boldsymbol x)},
\end{equation}
where $R(K)$ and $J$ act on each label with its own degree. These laws are
not consequences of the preceding assumptions on the actions: they are
properties of the correlations and are inputs to this section; they are
to be supplied, for the constructed correlations on the same hierarchy,
by the producer of the Euclidean field structure.
Section~\ref{sec:ope-estimates} uses all three. Haar measure
$dK$ on $O(4)$ is normalized to have mass one.

\subsection{Averaging on a common separated region}
For a fixed degree $n$, let $\mathcal U_n$ be the set of cluster configurations
$(\boldsymbol x,o)$ for which every exterior test occurring through degree
$n$ remains separated from the relative cluster after every orthogonal
rotation. This is an open set. Indeed, for a fixed finite test, separation
is an open condition in the configuration and the rotation. If it holds
for every rotation at a configuration, compactness of $O(4)$ supplies
finitely many rotation neighborhoods and a common configuration neighborhood
on which it still holds. Intersect over the finitely many tests.
The set is invariant under simultaneous translations, rotations and
permutations of the cluster. It contains every sufficiently small separated
cluster around its center, by the clearance argument of the preceding
section.

On $\mathcal U_n$ define
\begin{equation}\label{eq:ope-average}
 A_n(\boldsymbol x,\boldsymbol v;o)
   =\int_{O(4)}R_n(K)
       H_n(K^{-1}\boldsymbol x,R(K^{-1})\boldsymbol v;K^{-1}o)\,dK,
\end{equation}
where the action on each input label uses its own degree. Set $A_n=0$
outside $\mathcal U_n$. The integrand is continuous in $K$, takes values
in a finite-dimensional space, and is therefore integrable on the compact
group. Multilinearity passes through the integral; the zero extension
preserves it because the domain condition is independent of labels.

The map $A_n$ is smooth in positions on $\mathcal U_n$. Here one must not
assume differentiability of the representation in $K$ to justify
position derivatives. Expand the input labels and output in finite bases.
Each matrix entry of $R(K)$ is continuous and bounded on the compact
group. The remaining factors are derivatives in positions of separated
correlations evaluated at rotated coordinates. On a sufficiently small
closed configuration neighborhood times $O(4)$, all finitely many
separations have a positive minimum. The correlation derivatives of each
fixed order are therefore continuous and bounded on this compact set.
The chain rule contributes only bounded entries of the rotation matrix.
Thus each position derivative has a uniform integrable bound.
Differentiation under the integral, repeated for every finite order,
proves smoothness. The same finite-basis expansion justifies smooth
variation of labels when needed.

Left invariance of Haar measure gives rotation covariance. For a fixed
$L\in O(4)$ substitute $K=LH$ in
$A_n(L\boldsymbol x,R(L)\boldsymbol v;Lo)$. Since
$(LH)^{-1}L=H^{-1}$, the integrand becomes $R_n(L)$ times the integrand
in \eqref{eq:ope-average}. Hence
\[
 A_n(L\boldsymbol x,R(L)\boldsymbol v;Lo)
                         =R_n(L)A_n(\boldsymbol x,\boldsymbol v;o).
\]
This also holds outside the domain by its invariance. Translation and
permutation identities follow directly from the corresponding identities
for $H_n$ inside the integral. For a single field $v\in E_q$ at its
center, with $q\leq n$, \eqref{eq:ope-raw-single} shows that the integrand
is the constant
$R_n(K)\iota_{qn}R_q(K^{-1})v=\iota_{qn}v$. Therefore
$A_n((o),(v);o)=\iota_{qn}v$.

Now define the real-symmetric readout
\begin{equation}\label{eq:ope-real-average}
 C_n(\boldsymbol x,\boldsymbol v;o)
  =\tfrac12\bigl(A_n(\boldsymbol x,\boldsymbol v;o)
       +J_nA_n(\boldsymbol x,J\boldsymbol v;o)\bigr).
\end{equation}
It is still complex multilinear: an input scalar is conjugated by the
input $J$ and conjugated back by the output $J_n$. It is smooth on
$\mathcal U_n$, has the same translation, permutation and rotation
identities, and satisfies
\[
 C_n(\boldsymbol x,J\boldsymbol v;o)=J_nC_n(\boldsymbol x,\boldsymbol v;o),
 \qquad C_n((o),(v);o)=\iota_{qn}v\quad(q\leq n).
\]
The first identity uses $J^2=I$ and exchanges the two summands. The
second follows from the corresponding identity for $A_n$ and
compatibility of $J$ with inclusions.

\subsection{Compatible projections onto new degrees}
For every pair $m,q\geq0$, define a linear map
\[
 L_{mq}:E_q\to E_m,\qquad L_{mq}v=C_m((0),(v);0).
\]
The centered single-field configuration always belongs to $\mathcal U_m$.
The preceding identities give
\begin{equation}\label{eq:ope-single-compat}
 L_{mq}=\iota_{qm}\ (q\leq m),\qquad
 L_{mq}\iota_{rq}=L_{mr}\ (r\leq q).
\end{equation}
For the second equality, every exterior correlation used in the raw tests
is unchanged under the physical inclusion. This equality passes through
rotation averaging, since the inclusions commute with rotations, and
through conjugation averaging for the same reason. It holds even when
$m<r$; no extension of a low-degree inverse is being assumed.
The maps $L_{mq}$ intertwine rotations and conjugation as well.

Define recursively
\[
 K_m=I_{E_m}-\sum_{i<m}\iota_{im}K_iL_{im},\qquad
 P_{mq}=K_mL_{mq},\qquad
 V_m=\bigcap_{i<m}\ker P_{im}.
\]
The empty sum makes $K_0=I$. These are finite sums of fixed linear maps.
From \eqref{eq:ope-single-compat}, for $q\leq m$,
\begin{equation}\label{eq:ope-projection-recursion}
 P_{mq}v=\iota_{qm}v-\sum_{i<m}\iota_{im}P_{iq}v,
 \qquad
 \sum_{i\leq m}\iota_{im}P_{iq}v=\iota_{qm}v.
\end{equation}
Also $P_{mq}\iota_{rq}=P_{mr}$ whenever $r\leq q$.
If $q<m$, apply the second identity at degree $m-1$ and then include
in $E_m$. Its sum is the one subtracted in the first identity, so
\begin{equation}\label{eq:ope-higher-kills}
 P_{mq}=0\quad(q<m).
\end{equation}
For $v\in V_m$, the recursion with $q=m$ gives $P_{mm}v=v$.

\begin{lemma}\label{lem:ope-homogeneous-projections}
The range of $K_m$ is contained in $V_m$. If $v\in V_j$, then
$P_{mj}v=\delta_{mj}v$, interpreting the right side as zero when
$m\neq j$. Consequently
\[
 E_m=\bigoplus_{j\leq m}\iota_{jm}V_j.
\]
\end{lemma}
\begin{proof}
We prove the range assertion by induction on $m$. Fix $i<m$ and apply
$P_{im}$ to the defining formula for $K_mx$. Compatibility gives
\[
 P_{im}K_mx=P_{im}x-\sum_{j<m}P_{ij}(K_jL_{jm}x).
\]
By the induction hypothesis $K_jL_{jm}x\in V_j$. If $i<j$, its image
under $P_{ij}$ is zero by the definition of $V_j$; if $i>j$, it is zero
by \eqref{eq:ope-higher-kills}. At $i=j$, the identity on $V_i$
proved above makes the term $K_iL_{im}x=P_{im}x$. The sum therefore
cancels the first term. This proves $K_mx\in V_m$.
The asserted action on $V_j$ follows from the same three cases.

Take $q=m$ in \eqref{eq:ope-projection-recursion}. Each
$P_{jm}x=K_jL_{jm}x$ belongs to $V_j$, so this identity proves spanning.
If $\sum_{j\leq m}\iota_{jm}v_j=0$ with $v_j\in V_j$, apply
$P_{km}$ and use inclusion compatibility. All terms vanish except
$v_k$. Thus every $v_k=0$, proving directness.
\end{proof}

Induction in the formula for $K_m$ shows that it commutes with rotations
and conjugation: the inclusions and $L_{im}$ intertwine them, and the
previous $K_i$ do so by induction. It follows that $V_m$ is invariant
under both operations, since all its defining tests intertwine them.
No claim that the original arbitrary quotient lifts were covariant was
needed.

Define the homogeneous output at degree $m$ by
\[
 F_m(\boldsymbol x,\boldsymbol v;o)
                         =K_mC_m(\boldsymbol x,\boldsymbol v;o)\in V_m.
\]
It is multilinear, has the same symmetries as $C_m$, and is smooth on
$\mathcal U_m$. At a single centered field $v\in E_q$ it equals
$P_{mq}v$. In particular the finite packet
\[
 \mathcal P_D(\boldsymbol x,\boldsymbol v;o)
          =\sum_{m\leq D}\iota_{mD}F_m(\boldsymbol x,\boldsymbol v;o)
\]
satisfies the exact compatibility
\begin{equation}\label{eq:ope-packet-projection}
 P_{jD}\mathcal P_D=F_j\qquad(j\leq D),
\end{equation}
by Lemma~\ref{lem:ope-homogeneous-projections}. On one centered field of
degree $q\leq D$, the packet equals $\iota_{qD}v$ by
\eqref{eq:ope-projection-recursion}.

\subsection{Real channel bases and scalar coefficients}
The conjugation $J_m$ restricts to an antilinear involution of $V_m$.
Its fixed real vector space has a real basis $(e_{m,\alpha})$; this is
also a complex basis of $V_m$. Indeed every $v\in V_m$ decomposes as
\[
 v=\frac{v+J_mv}{2}
       +i\frac{v-J_mv}{2i},
\]
and both displayed components are fixed by $J_m$. This proves complex
spanning. If $\sum_\alpha z_\alpha e_{m,\alpha}=0$, applying $J_m$
and then adding and subtracting the two equations shows separately that
the real and imaginary coefficient sums vanish. Real linear independence
forces every $z_\alpha=0$.

Choose these bases once for every $m$ and write
\begin{equation}\label{eq:ope-final-coefficients}
 F_m(\boldsymbol x,\boldsymbol v;o)
   =\sum_\alpha c_{\boldsymbol v}^{m,\alpha}(\boldsymbol x;o)
                                                   e_{m,\alpha}.
\end{equation}
The coordinate functionals are continuous because the spaces are finite
dimensional. Thus the coefficients inherit smoothness, multilinearity,
translation and permutation symmetry; reality gives conjugation of their
values under conjugation of all input labels. Rotation covariance holds
for the field-valued output, or equivalently through the finite matrix of
$R_m$ on this basis. The individual basis vectors need not be rotation
invariant.

Let $\ell_{\chi,m,o,\boldsymbol y,\boldsymbol w}$ be the linear
functional obtained by inserting an $E_m$ field at $o$ alongside separated
spectators, in any phase $\chi$ or in the mixed state. Physical inclusion
compatibility gives
$\ell_{\chi,D}\iota_{mD}=\ell_{\chi,m}$. Therefore evaluation of the
packet is exactly
\[
 \ell_{\chi,D}(\mathcal P_D)
  =\sum_{m\leq D}\sum_\alpha
    c_{\boldsymbol v}^{m,\alpha}(\boldsymbol x;o)
    S_\chi(e_{m,\alpha},\boldsymbol w;o,\boldsymbol y).
\]
This identifies the final scalar coefficients with one fixed covariant
field prescription, independent of phase, spectators and truncation.
It does not yet estimate the difference from the original cluster
correlation. That quantitative difference is the next proof obligation.
